\subsection{Evaluation metrics}
\label{sec:scoring}

\phantomsection
\label{sec:scoring-why}
The recorded sequence provides one solution, but other sequences may produce the required
construction. We therefore assess correctness by execution rather than agreement with the
reference actions. Distance between the constructed crease pattern and the input crease pattern
is a separate measure of geometric progress; it does not compare action sequences or replace
final-state verification.

\subsubsection{Verification}
\label{sec:scoring-what}

We execute the proposed sequence and compare its accumulated creases with the input pattern,
including mountain--valley assignments. We then compare the final stack with the target after
alignment by a single global translation, rotation, or turnover. A turnover reflects the geometry,
reverses the stack order, and flips face orientations. The same transformation must align all
layers; independent layer transformations and arbitrary layer permutations are not accepted.
A successful construction must satisfy both crease-pattern and folded-state agreement.

Neither comparison is tolerance-free. The crease comparator and the shared folded-state
comparator both match positions within $2\times10^{-6}$; Table~\ref{tab:tolerances} lists every
tolerance in the pipeline, and Appendix~\ref{app:comparison-details} describes the alignment and
representation details. These final checks are distinct from the three optional
levels of target feedback supplied during an episode.

\subsubsection{Reported metrics}
\label{sec:scoring-metrics}

We report solve rate by fold-count group, with the number of attempts and the applicable run
budgets. Termination reasons distinguish successful submission, unsuccessful completion, budget
limits, and repetition-related stopping. Query counts describe search effort; for unsuccessful
attempts they are consumed queries, not observed queries to solution.

% EDITING NOTE: Compute CP edit distance for the final valid state of every saved attempt,
% including unsuccessful and interrupted runs. Define edits, costs, normalization, and treatment
% of missing state artifacts before adding numerical results. See IMPORTANT_TODO.md.
We measure CP edit distance between the accumulated crease pattern in the last recoverable
state and the supplied pattern, in original-sheet coordinates. Collinear mountain and valley
intervals are merged separately, excluding sheet boundaries. If $m$ intervals match under the
crease comparator, the insertion/deletion distance is
$d_{\mathrm{CP}}=|C_{\mathrm{target}}|+|C_{\mathrm{candidate}}|-2m$.
Each missing or extra interval costs one; changing an assignment costs two. Partially overlapping
intervals do not match unless their endpoints agree within tolerance. We also report
$d_{\mathrm{CP}}/|C_{\mathrm{target}}|$, which can exceed one. Normalized distance is undefined
for an empty target. Means weight attempts equally and exclude missing state artifacts,
whose counts are reported separately. This metric compares crease patterns, not fold sequences,
and zero distance does not establish folded-state agreement.

\phantomsection
\label{sec:scoring-validate}
The comparator that scores every reported run is validated by a test that must accept the global
transformations and reject everything else. Acceptance is checked across rotations, turnovers and
translations, in both comparison directions, and under changes to polygon start vertex and winding
that do not change the folded object. Rejection is checked for a turnover without stack reversal,
a turnover without parity flips, a coordinate mirror presented as a turnover, a permuted layer
order, a missing layer, a uniform scaling, a per-layer transformation that no single isometry
explains, and an empty stack. The polygons used are asymmetric, so an incorrect reflection cannot
pass through a symmetry of the square. This validates the implementation on the scoring path. The
corpus-side comparator is a separate implementation with its own negative control, described in
Appendix~\ref{app:comparison-details}, and does not stand in for it.
